% !TEX root = ../Thesis.tex
% ============================================================
\chapter{Implementation} \label{chapter:implementation}
% ============================================================

% EDITORIAL NOTE. This chapter exists to make two claims checkable: that the PCST
% implementation used here behaves exactly like the published one, and that every
% number in Chapter 6 can be regenerated. Anything that does not serve one of
% those does not belong here.

The software is a Python package, \texttt{ikgqa}, organised so that each concern
in Chapter~\ref{chapter:methodology} lives in one place and can be exercised on
its own. This chapter describes the parts whose correctness the results depend
on, and how each is established.

\section{Structure} \label{sec:implementation-structure}

\begin{table}[t]
\centering
\caption{Modules and their responsibilities.}
\label{tab:modules}
\begin{tabular}{ll}
\toprule
Module & Responsibility \\
\midrule
\texttt{ikgqa.graph} & a graph with text and embeddings, validated on construction \\
\texttt{ikgqa.encoders} & bag-of-words and sentence-transformer encoders \\
\texttt{ikgqa.pcst} & the PCST algorithm and a verbatim copy of the published one \\
\texttt{ikgqa.retrieval} & the five strategies, plus two-stage candidate generation \\
\texttt{ikgqa.eval} & the metric, the parameter sweep, and curve extraction \\
\texttt{ikgqa.data} & the Neo4j loader, terminology resolution, the replica \\
\bottomrule
\end{tabular}
\end{table}

Two design decisions in Table~\ref{tab:modules} are load-bearing rather than
stylistic.

The graph type validates its inputs when constructed: one embedding row per node,
one per edge, matching dimensions, and edge endpoints within range. The reason is
that none of these mismatches crashes. A graph with the wrong number of embedding
rows produces a plausible but wrong subgraph, and a research bug of that kind can
survive for months. Validating once converts it into an error message.

Every retrieval strategy implements one interface -- given a graph and an embedded
question, return node and edge identifiers -- and nothing else. The strategies do
not embed text, do not score themselves, and do not call a model. That narrowness
is what allows the evaluation harness to treat them as interchangeable, which is
the premise of the entire comparison.

\section{Fidelity to the Published Implementation} \label{sec:implementation-fidelity}

Chapter~\ref{chapter:results} attributes behaviour to G-Retriever's design. That
attribution is only legitimate if the implementation used here is G-Retriever's,
so fidelity is established rather than asserted.

A verbatim copy of the published retrieval routine is kept in the package,
confirmed byte-identical to the upstream source, and is never edited. It serves
as a test oracle: the instrumented implementation used throughout this thesis is
required to return the same node set, the same edge set and the same serialised
description as the copy, on curated graphs chosen to exercise each branch and on
randomly generated graphs. Any divergence is a test failure rather than a
judgement call.

Three deviations from upstream exist, and all three are deliberate, localised and
documented at the point of change.

\paragraph{Edge identifier ordering.} Decoding concatenates the solver's real
edges with the edges recovered from virtual nodes and never re-sorts, so the
returned sequence is in solver order. This is harmless for a set-based metric and
silently misleading for anything comparing sequences positionally. The ordering is
normalised in the interface layer, leaving the algorithm byte-faithful.

\paragraph{Tie-breaking.} Selecting the top $k$ nodes requires a total order, and
where similarities are exactly equal the published implementation inherits the
order of its tensor library's selection routine. Because ties are the normal case
on this graph rather than an edge case, the tie-breaking rule is made an explicit
parameter: one setting reproduces the published behaviour exactly, another is
deterministic without that dependency. Every experiment records which was used.

\paragraph{A pinned numerical dependency.} The PCST solver package is compiled
against the NumPy 1.x binary interface. Under NumPy 2 it imports without error and
returns arrays of the correct shape filled with meaningless values. The dependency
is pinned, and a runtime check verifies the solver against a known-answer instance
before any experiment, because a silent numerical failure is the one class of bug
that testing after the fact cannot catch.

\section{The Loader} \label{sec:implementation-loader}

The loader turns one patient into the package's graph type using read-only
Cypher. Folding of shared terminology nodes is performed in the query rather than
in Python, so that only flat rows cross the network and no unnecessary node is
ever materialised. Patients are addressed by position in an ordered result, which
keeps pseudo-identifiers on the server.

Fetching and assembly are separate functions. Assembly is pure -- rows in, graph
out -- so it is unit tested against fixture rows with no database, while the
Cypher itself is validated by running it. This separation is what allowed the
assembly logic to be developed and tested before database access was available,
and it is what lets the replica of Section~\ref{sec:implementation-replica} reuse
the real assembly path.

Shared clinical entities are deduplicated during assembly: repeated measurements
of one analyte reference a single analyte node, and repeated administrations of
one drug a single drug node. The resulting graph is deliberately not a star
around the patient, because a star would make connectivity vacuous and leave the
selection problem with nothing to decide.

\section{The Replica} \label{sec:implementation-replica}

Access to the clinical graph is intermittent, and every component downstream of
the loader needs a graph to run against. The replica generates a synthetic
patient whose structure matches a measured one exactly.

Fidelity is achieved by construction rather than by imitation: the generator
emits rows in the shape the Cypher returns and passes them to the same assembly
function the loader uses. The structure is therefore identical by definition, and
the assembly code is exercised by every offline experiment rather than only when
the database is reachable. Achieved counts are printed against measured targets
on every run, so the one quantity that does not match exactly -- the number of
distinct texts, eleven per cent high -- is visible rather than buried.

The limits of the replica are stated in Section~\ref{sec:results-scope} and
repeated in the module itself, because a synthetic artefact that is faithful in
one respect invites being trusted in others.

\section{Testing} \label{sec:implementation-testing}

The suite covers the algorithm's equivalence to the published implementation, the
metric's behaviour on undefined cases, the loader's assembly against fixtures,
the terminology tiers, and the two-stage index mapping.

Two areas receive disproportionate attention because their failures are silent.
The first is the retrieval contract -- identifiers sorted, unique, in range, and
the node set closed under the returned edges -- asserted after every strategy in
every test, which is how the edge-ordering deviation above was found. The second
is the translation of identifiers between a narrowed subgraph and the original
graph: a result returned in the wrong index space consists of plausible
identifiers, computes a plausible recall, and is entirely wrong. It is tested
directly, with a generator and an inner retriever chosen so the two index spaces
cannot coincide by accident, rather than only through end-to-end recall.

\section{Reproducibility} \label{sec:implementation-reproducibility}

Every experiment is a script that writes tidy result files, and every figure and
table in Chapter~\ref{chapter:results} names the script that produced it.
Randomness is seeded, tie-breaking is an explicit parameter, and the offline
encoder is deterministic given its vocabulary, so a rerun reproduces a run.
Credentials are read from the environment and never appear in code or in results.
